\documentclass[10pt,a4paper]{amsart}
\usepackage[margin=19mm]{geometry}
\usepackage{amsmath,amssymb,mathtools}
\usepackage[scr=rsfs,cal=euler]{mathalfa}
\usepackage{xfrac}
\usepackage[unicode,hidelinks,bookmarks=false]{hyperref}
\newcommand{\ie}{i.e.\ }
\newcommand{\cf}{cf.\ }
\newcommand{\resp}{respectively\ }
\newcommand{\Subsection}{\subsection}
\providecommand{\nobreakdash}{\nobreak\hbox{-}}
\setlength{\emergencystretch}{2em}
\multlinegap0pt
\usepackage{xcolor}
\usepackage{bm}
\usepackage{mathtools}
\usepackage{amssymb}
\usepackage{tcolorbox}
\usepackage{float}
\usepackage{enumerate}
\usepackage{dsfont}
\newtheorem{lemma}{Lemma}
\def\RR{\mathbb R}
\title{Learning to Detect Cyber Attacks: Neural Anomaly Detection for Cybersecurity with Theoretical Insights}
\author{Tian-Yi Zhou$^*$ , Matthew Lau$^*$ , Jizhou Chen , Wenke Lee , Xiaoming Huo}
\date{Source: arXiv:2409.08521v2 (CC BY 4.0)}
\begin{document}
\maketitle

\noindent\emph{Source and licence.} Learning to Detect Cyber Attacks: Neural Anomaly Detection for Cybersecurity with Theoretical Insights. Version arXiv:2409.08521v2 (CC BY 4.0), distributed by arXiv under the Creative Commons Attribution 4.0 International licence, http://creativecommons.org/licenses/by/4.0/, as recorded in the arXiv licence metadata for this exact version and shown on the abstract page https://arxiv.org/abs/2409.08521v2. Full licence notice: \texttt{licenses/arxiv-2409.08521-CC-BY-4.0.txt}.\par\smallskip

\noindent\textit{Editorial scope.} Counted unit: the article's lemma lemma:estimationII (source0.tex lines 2340--2352). The article's complete original proof of it is delivered with it (source0.tex lines 2354--2496, one \texttt{proof} environment of 143 lines). No mathematics is written, repaired or re-derived: the statement and the proof are the article's own lines, in the article's own notation, and no retained line is substituted or re-flowed.  Retained uncounted as the source-local context the proof needs: lines 2129--2131; lines 2188--2190; lines 2294--2307.  Nothing was omitted from any retained span, so the package carries no omission record.  No retained line contains a citation, so the wrapper carries no bibliography and no bibliography entry.  No retained line contains a figure environment, an image-inclusion command or drawing code, so no image is delivered and the package declares no asset dependency.  Editorial formatting only, in addition: an AMS \texttt{amsart} preamble that restates the article's own declarations and macros used by the retained lines, namely the environments \texttt{lemma} on separate counters, together with the standard packages those lines need.  The retained environments print in this document as Lemma 1 in order of appearance, whereas the article prints the same items with the numbering of its own full document.  The complete native source of the article as distributed in the arXiv e-print for this exact version is bundled unchanged as \texttt{sources/165185/source0.tex}.
 \begin{equation}\label{Young}
        a\cdot b \leq \frac{a^p}{p}+ \frac{b^{p^*}}{p^*} \qquad \text{with    } a\geq 0,b \geq 0, p>1,p^*>1 \text{ and } \frac{1}{p}+ \frac{1}{p^*} = 1.
    \end{equation}

    \begin{equation}\label{coveringeqn}
        \log \mathcal{N}(\epsilon,\mathcal{H}_\tau) \leq c_{\alpha, d} m^2  N \log ((\tau\epsilon)^{-1}mN),
    \end{equation}

\begin{lemma}\label{uniformBernstein}
Let ${\mathcal H}$ be a set of measurable functions from ${\mathcal X}$ to $[-1, 1]$. Suppose the noise condition (Assumption 2.1) holds for some $q\in [0, \infty]$ and $c_0 >0$. Then for every
$\epsilon>0$, 
$$\sup_{f\in {\mathcal H}} \frac{\varepsilon (f) - \varepsilon (f_c) - \left(\varepsilon_{T, T'} (f) - \varepsilon_{T, T'} (f_c)\right)}{\left(\left(\varepsilon (f) - \varepsilon (f_c)\right)^{\frac{q}{q+1}} + \epsilon^{\frac{q}{q+1}}\right)^{1/2}} \leq 5 \epsilon^{1-{\frac{q}{2(q+1)}}}
$$
with probability at least 
\begin{eqnarray*}
1-  {\mathcal N}\left(\epsilon, {\mathcal H}\right) 
\left\{\exp \left\{-
{n \epsilon^{2-{\frac{q}{q+1}}} \over s\left(10 (c_0)^{\frac{1}{q+1}} +2 \epsilon^{1-{\frac{q}{q+1}}}\right)}\right\}
 +\exp \left\{-
{n' \epsilon^{2-{\frac{q}{q+1}}} \over (1-s) \left(10 (c_0)^{\frac{1}{q+1}} +2 \epsilon^{1-{\frac{q}{q+1}}}\right)}\right\}\right\}. 
\end{eqnarray*}
\end{lemma}

\begin{lemma}\label{lemma:estimationII}
     Let $0\leq q \leq \infty$. Let $\alpha,r >0$, and integers $m \geq 1$ and \\$N \geq \max \left\{(\alpha +1)^d, (r+1)e^d\right\}$.
    Suppose the noise condition (Assumption 2.1) holds for some $q$ and constant $c_0>0$. Suppose $\tau \geq \max\left\{\frac{1}{n}, \frac{1}{n^\prime}\right\}$. 
    For any $0 <\delta < 1$ and $n\geq 3$, with probability $1-\frac{\delta}{2}$, we have 
\begin{eqnarray*}
  &&   \varepsilon (\Hat{f}_{T, T^\prime,\phi}) - \varepsilon (f_c) - \left(\varepsilon_{T, T'} (f) - \varepsilon_{T, T'} (f_c)\right)\\
 &\leq&  C_{q,\alpha, d} \left(\log \left(\frac{4}{\delta}\right) + m^2 N (\log(nmN)+\log(n^\prime mN))\right)^{\frac{q+2}{q+1}} \\
 &&\quad\quad\cdot\left(  \max\left\{  \frac{s}{n}
\log \left(\frac{n}{s}\right), \frac{(1-s)}{n^\prime}\log \left(\frac{n^\prime}{1-s}\right) \right\} \right)^{\frac{q+2}{q+1}} 
+ \frac{\varepsilon (\Hat{f}_{T, T^\prime,\phi}) - \varepsilon (f_c)}{2},
\end{eqnarray*}
where $C_{q,\alpha, d}$ is a positive constant depending only on $q,c_0, \alpha$, and $d$. 
\end{lemma}

\begin{proof}[Proof of Lemma \ref{lemma:estimationII}]
    Lemma \ref{uniformBernstein} tells us that, for every $0<\epsilon \leq 1$, with probability at least 
    \begin{eqnarray*}
&&1-  {\mathcal N}\left(\epsilon, {\mathcal H_\tau}\right) \left\{\exp \left\{-
{n \epsilon^{2-{\frac{q}{q+1}}} \over s\left(10 (c_0)^{\frac{1}{q+1}} +2 \epsilon^{1-{\frac{q}{q+1}}}\right)}\right\}
+\exp \left\{-
{n' \epsilon^{2-{\frac{q}{q+1}}} \over (1-s) \left(10 (c_0)^{\frac{1}{q+1}} +2 \epsilon^{1-{\frac{q}{q+1}}}\right)}\right\}\right\}\\
&\geq& 1-  {\mathcal N}\left(\epsilon, {\mathcal H_\tau}\right) \left\{\exp \left\{-
{n \epsilon^{2-{\frac{q}{q+1}}} \over s\left(10 (c_0)^{\frac{1}{q+1}} +2 \right)}\right\} +\exp \left\{-
{n' \epsilon^{2-{\frac{q}{q+1}}} \over (1-s) \left(10 (c_0)^{\frac{1}{q+1}} +2 \right)}\right\}\right\}\\
&\stackrel{\eqref{coveringeqn}}{\geq}& 1- \exp \left\{c_{\alpha, d} m^2 N \log ((\tau\epsilon)^{-1}mN) -
{n \epsilon^{2-{\frac{q}{q+1}}} \over s\left(10 (c_0)^{\frac{1}{q+1}} +2 \right)}\right\}\\
&& -\exp \left\{c_{\alpha, d} m^2 N \log ((\tau\epsilon)^{-1}mN)-
{n' \epsilon^{2-{\frac{q}{q+1}}} \over (1-s) \left(10 (c_0)^{\frac{1}{q+1}} +2 \right)}\right\} ,
\end{eqnarray*} 
there holds
$$\sup_{f\in {\mathcal H}} \frac{\varepsilon (f) - \varepsilon (f_c) - \left(\varepsilon_{T, T'} (f) - \varepsilon_{T, T'} (f_c)\right)}{\left(\left(\varepsilon (f) - \varepsilon (f_c)\right)^{\frac{q}{q+1}} + \epsilon^{\frac{q}{q+1}}\right)^{1/2}} \leq 5 \epsilon^{1-{\frac{q}{2(q+1)}}},
$$
which implies 
\begin{equation}\label{5epsilon}
    \varepsilon (f) - \varepsilon (f_c) - \left(\varepsilon_{T, T'} (f) - \varepsilon_{T, T'} (f_c)\right)\leq 5 \epsilon^{1-{\frac{q}{2(q+1)}}}\left(\left(\varepsilon (f) - \varepsilon (f_c)\right)^{\frac{q}{q+1}} + \epsilon^{\frac{q}{q+1}}\right)^{1/2}, \quad \forall f\in \mathcal{H}_\tau. 
\end{equation}

We choose $\tau$ satisfying $\tau \geq \max\{\frac{1}{n}, \frac{1}{n^\prime}\}$.
Then we know $\tau^{-1} \leq n$ and $\tau^{-1} \leq n^\prime$.
We set the above confidence bound to be at least $1-\delta/2$. We need to find $\epsilon$ that satisfies the following two inequalities:
$$\exp \left\{c_{\alpha, d} m^2 N \log (\epsilon^{-1}nmN) -
{n \epsilon^{2-{\frac{q}{q+1}}} \over s\left(10 (c_0)^{\frac{1}{q+1}} +2 \right)}\right\} \leq \frac{\delta}{4}$$
and 
$$\exp \left\{c_{\alpha, d} m^2 N \log (\epsilon^{-1}n^\prime mN)-
{n' \epsilon^{2-{\frac{q}{q+1}}} \over (1-s) \left(10 (c_0)^{\frac{1}{q+1}} +2 \right)}\right\} \leq \frac{\delta}{4}.$$ 

Let $\hat \epsilon = \epsilon^{2-{\frac{q}{q+1}}}$ and let $c_q = 10 (c_0)^{\frac{1}{q+1}} +2$. We see that $c_q$ is a positive constant depending only on $q$ and $c_0$. We need to find $\hat \epsilon>0$ such that 
\begin{equation}\label{firstinq}
    \frac{c_{\alpha, d}}{2-{\frac{q}{q+1}}} m^2 N \log (\hat \epsilon^{-1}) - \frac{n\hat \epsilon }{s c_{q}}  \leq \log \left(\frac{\delta}{4}\right) -  c_{\alpha, d} m^2 N \log(nmN)    
\end{equation}
and 
\begin{equation}\label{secondinq}
    \frac{c_{\alpha, d}}{2-{\frac{q}{q+1}}} m^2 N \log (\hat \epsilon^{-1}) - \frac{n^\prime \hat \epsilon }{(1-s) c_{q}}  \leq \log \left(\frac{\delta}{4}\right) -  c_{\alpha, d} m^2 N \log(n^\prime mN).    
\end{equation}
We define a function $T:(0, \infty) \to \RR$ given by 
\begin{equation}\label{Tfunction}
     T(x) = \frac{c_{\alpha, d}}{2-{\frac{q}{q+1}}} m^2 N \log (x^{-1}) - \frac{nx }{sc_{q}}.
\end{equation}     
We notice that $T$ is a decreasing function. Then, we take
\begin{equation}\label{B}
    B:=c_{q} \left(\frac{2c_{\alpha, d}}{2-{\frac{q}{q+1}}} m^2 N + \log \left(\frac{4}{\delta}\right) + c_{\alpha, d} m^2 N\log(nmN)\right).  
\end{equation}
For $n \geq 3$ (i.e., $\log (n) >1$), there holds 
\begin{equation}\label{Bs}
    \frac{B s\log \left(\frac{n}{s}\right)}{n}
 \geq \frac{s}{n}.
\end{equation}
It follows that  
\begin{eqnarray*}
   && T\left( \frac{Bs\log \left(\frac{n}{s}\right)}{n}
\right)\\
   &\stackrel{\eqref{Tfunction}}{=}& \frac{c_{\alpha, d}}{2-{\frac{q}{q+1}}} m^2 N \log \left(\frac{n}{Bs\log (\frac{n}{s})}\right) - \frac{n}{sc_{q}}\left(\frac{Bs\log (\frac{n}{s})}{n}\right)\\
   &\stackrel{\eqref{B}, \eqref{Bs}}{\leq}& \frac{c_{\alpha, d}}{2-{\frac{q}{q+1}}} m^2 N \log \left(\frac{n}{s}\right)\\
   &&\qquad- \log \left(\frac{n}{s}\right)\left(\frac{c_{\alpha, d}}{2-{\frac{q}{q+1}}} m^2 N + \log \left(\frac{4}{\delta}\right) + c_{\alpha, d} m^2 N \log(nmN)\right)\\
   &=&  - \log \left(\frac{n}{s}\right) \left(\log \left(\frac{4}{\delta}\right) + c_{\alpha, d} m^2 N \log(nmN)\right)\\
&\stackrel{\because \log (\frac{n}{s}) > 1}{\leq}& \log \left(\frac{\delta}{4}\right) -  c_{\alpha, d} m^2  N\log(nmN). 
\end{eqnarray*}
Since $T$ is a decreasing function, we know that if
\begin{eqnarray*}
  \hat \epsilon \geq  \frac{Bs\log \left(\frac{n}{s}\right)}{n}
,
\end{eqnarray*}
then such a $\hat \epsilon$ satisfies the first inequality \eqref{firstinq}.
In the same way, we know that if
\begin{eqnarray*}
    \hat\epsilon \geq \frac{B^\prime(1-s)}{n^\prime}\log \left(\frac{n^\prime}{1-s}\right)
\end{eqnarray*}
with \begin{eqnarray*}
    B^\prime:=c_{q} \left(\frac{2c_{\alpha, d}}{2-{\frac{q}{q+1}}} m^2 N + \log \left(\frac{4}{\delta}\right) + c_{\alpha, d} m^2 N\log(n^\prime mN)\right),  
\end{eqnarray*}
then such a $\hat\epsilon$ satisfies the second inequality \eqref{secondinq}.

We choose $\hat \epsilon$ to be  
\begin{eqnarray*}
   \hat\epsilon =  \max \left\{  \frac{Bs\log \left(\frac{n}{s}\right)}{n}
, \frac{B^\prime(1-s)}{n^\prime}\log \left(\frac{n^\prime}{1-s}\right) \right\}.
\end{eqnarray*}
Thus, $\hat\epsilon $ satisfies both inequalities \eqref{firstinq} and \eqref{secondinq}. Then, we have  \begin{equation}\label{choice_of_epsilon}
    \epsilon =   \left(\max\left\{  \frac{Bs\log \left(\frac{n}{s}\right)}{n}
, \frac{B^\prime(1-s)}{n^\prime}\log \left(\frac{n^\prime}{1-s}\right) \right\}\right)^{\frac{q+1}{q+2} }
\end{equation}
because $\hat \epsilon = \epsilon^{2-{\frac{q}{q+1}}}$.

Now, we take $f = \Hat{f}_{T, T^\prime,\phi} \in \mathcal{H}_\tau$ and the above choice of $\epsilon$ \eqref{choice_of_epsilon}, we have,  with probability at least $1-\delta/2$, 
\begin{eqnarray*}
&&  \varepsilon (\Hat{f}_{T, T^\prime,\phi}) - \varepsilon (f_c) - \left(\varepsilon_{T, T'} (f) - \varepsilon_{T, T'} (f_c)\right)\\
  &\stackrel{\eqref{5epsilon}}{\leq}& 5 \epsilon^{1-{\frac{q}{2(q+1)}}}\left(\left(\varepsilon (\Hat{f}_{T, T^\prime,\phi}) - \varepsilon (f_c)\right)^{\frac{q}{q+1}} + \epsilon^{\frac{q}{q+1}}\right)^{1/2}\\
  &\stackrel{\eqref{choice_of_epsilon}}{=}& 5 \left(  \max\left\{ B \frac{s}{n}
\log \left(\frac{n}{s}\right),B^\prime \frac{(1-s)}{n^\prime}\log \left(\frac{n^\prime}{1-s}\right) \right\}\right)^{1/2}\\
&&\left(\left(\varepsilon (\Hat{f}_{T, T^\prime,\phi}) - \varepsilon (f_c)\right)^{\frac{q}{q+1}} + \left( \max\left\{ B  \frac{s}{n}
\log \left(\frac{n}{s}\right), B^\prime \frac{(1-s)}{n^\prime}\log \left(\frac{n^\prime}{1-s}\right) \right\}\right)^{\frac{q}{q+2} } \right)^{1/2} \\
&\leq& 5 \left(  \max\left\{ B \frac{s}{n}
\log \left(\frac{n}{s}\right), B^\prime \frac{(1-s)}{n^\prime}\log \left(\frac{n^\prime}{1-s}\right) \right\} \right)^{1/2} \left(\varepsilon (\Hat{f}_{T, T^\prime,\phi}) - \varepsilon (f_c)\right)^{\frac{q}{2(q+1)}} \\
&& + 5 \left(  \max\left\{ B \frac{s}{n}
\log \left(\frac{n}{s}\right), B^\prime\frac{(1-s)}{n^\prime}\log \left(\frac{n^\prime}{1-s}\right) \right\}\right)^{\frac{q+1}{q+2}}.
\end{eqnarray*}
Here, we have used $\sqrt{a+b} \leq \sqrt{a} + \sqrt{b}$ for all $a,b >0$ in the last inequality.  To further tighten this upper bound, we apply Young's Inequality \eqref{Young} with\\
$a=5 \left(  \max\left\{  B\frac{s}{n}
\log \left(\frac{n}{s}\right),  B^\prime \frac{(1-s)}{n^\prime}\log \left(\frac{n^\prime}{1-s}\right) \right\} \right)^{1/2}$,
$b= \left(\varepsilon (\Hat{f}_{T, T^\prime,\phi}) - \varepsilon (f_c)\right)^{\frac{q}{2(q+1)}}, p=\frac{2(q+1)}{q+2}$, $ p^*=\frac{2(q+1)}{q}$ and get
\begin{eqnarray}\label{apply_Young}
  &&  5 \left(  \max\left\{  B\frac{s}{n}
\log \left(\frac{n}{s}\right),  B^\prime \frac{(1-s)}{n^\prime}\log \left(\frac{n^\prime}{1-s}\right) \right\} \right)^{1/2} \left(\varepsilon (\Hat{f}_{T, T^\prime,\phi}) - \varepsilon (f_c)\right)^{\frac{q}{2(q+1)}} \nonumber\\
&\leq&   \frac{\left( 5\left(  \max\left\{ B \frac{s}{n}
\log \left(\frac{n}{s}\right),B^\prime \frac{(1-s)}{n^\prime}\log \left(\frac{n^\prime}{1-s}\right) \right\} \right)^{1/2}\right)^{\frac{2(q+1)}{q+2}}}{\frac{2(q+1)}{q+2}} + \frac{\varepsilon (\Hat{f}_{T, T^\prime,\phi}) - \varepsilon (f_c)}{\frac{2(q+1)}{q}} \nonumber\\
&\leq&5^{\frac{2(q+1)}{q+2}} \left(  \max\left\{  B\frac{s}{n}
\log \left(\frac{n}{s}\right), B^\prime \frac{(1-s)}{n^\prime}\log \left(\frac{n^\prime}{1-s}\right) \right\} \right)^{\frac{q+1}{q+2}} + \frac{\varepsilon (\Hat{f}_{T, T^\prime,\phi}) - \varepsilon (f_c)}{2}.
\end{eqnarray}
Furthermore, we have 
\begin{eqnarray}\label{Bbound}
    B&\stackrel{\eqref{B}}{=}&c_{q} \left(\frac{2c_{\alpha, d}}{2-{\frac{q}{q+1}}} m^2 N + \log \left(\frac{4}{\delta}\right) + c_{\alpha, d} m^2 N\log(nmN)\right) \nonumber\\
    &\leq& c_{q} \left(\frac{2c_{\alpha, d}}{2-{\frac{q}{q+1}}} + c_{\alpha, d}\right)\left(m^2 N + \log \left(\frac{4}{\delta}\right) + m^2 N\log(nmN)\right)\nonumber\\
     &\leq&  c_{q,\alpha,d}\left(\log \left(\frac{4}{\delta}\right) + m^2 N (1+\log(nmN))\right),
\end{eqnarray}
where $c_{q,\alpha, d} = c_{q} \left(\frac{2c_{\alpha, d}}{2-{\frac{q}{q+1}}} + c_{\alpha, d}\right) \geq 1$ is a positive constant depending only on $q, c_0 , \alpha$, and $d$. 
Similarly, we know 
\begin{equation}\label{Bprimebound}
    B^\prime \leq  c_{q,\alpha,d}\left(\log \left(\frac{4}{\delta}\right) + m^2 N (1+\log(n^\prime mN))\right). 
\end{equation}
 
Finally, we get with probability at least $1-\delta/2$,
\begin{eqnarray*}
    &&  \varepsilon (\Hat{f}_{T, T^\prime,\phi}) - \varepsilon (f_c) - \left(\varepsilon_{T, T'} (f) - \varepsilon_{T, T'} (f_c)\right)\\ 
    &\stackrel{\eqref{apply_Young}}{\leq}&  5^{\frac{2(q+1)}{q+2}}  \left(  \max\left\{ B \frac{s}{n}
    \log \left(\frac{n}{s}\right), B^\prime \frac{(1-s)}{n^\prime}\log \left(\frac{n^\prime}{1-s}\right) \right\} \right)^{\frac{q+1}{q+2}} + \frac{\varepsilon (\Hat{f}_{T, T^\prime,\phi}) - \varepsilon (f_c)}{2} \\
 &\leq&  5^{\frac{2(q+1)}{q+2}}  \left(\max\{B, B^\prime\}\right)^{\frac{q+1}{q+2}}\left(  \max\left\{  \frac{s}{n}
\log \left(\frac{n}{s}\right), \frac{(1-s)}{n^\prime}\log \left(\frac{n^\prime}{1-s}\right) \right\} \right)^{\frac{q+1}{q+2}}\\
&&\qquad\qquad+ \frac{\varepsilon (\Hat{f}_{T, T^\prime,\phi}) - \varepsilon (f_c)}{2} \\
&\stackrel{\eqref{Bbound}, \eqref{Bprimebound}}{\leq}& C_{q,\alpha, d} \left(\log \left(\frac{4}{\delta}\right) + m^2 N (\log(nmN)+\log(n^\prime mN))\right)^{\frac{q+1}{q+2}} \\
&& \cdot \left(  \max\left\{  \frac{s}{n}
\log \left(\frac{n}{s}\right), \frac{(1-s)}{n^\prime}\log \left(\frac{n^\prime}{1-s}\right) \right\} \right)^{\frac{q+1}{q+2}} 
 + \frac{\varepsilon (\Hat{f}_{T, T^\prime,\phi}) - \varepsilon (f_c)}{2}, 
\end{eqnarray*}
where $C_{q,\alpha, d}$ is a positive constant depending only on $q,c_0, \alpha$, and $d$. 

The proof is complete.
\end{proof}

\end{document}
